% GENERATED FILE. Edit canonical lessons, then run npm run generate:books.
% canonical-source-hash: fnv1a64:3a5d3e108cc7f46f
% canonical-lessons: IT-C413-tamponare, IT-C413-tamponamento, IT-C413-ingorgo, IT-C413-deviazione, IT-C413-deviare

\chapter{Incidenti e code}
\label{ch:it-incidenti-e-code}
\hlchaptermodality{413}

\begin{chapteropening}
\textbf{By the end of this chapter:} \emph{I can report trouble on the road: a bump from behind, a pile-up, a traffic jam and a diversion.}

Complete the last lesson of chapter 413: I can report trouble on the road: a bump from behind, a pile-up, a traffic jam and a diversion.
\end{chapteropening}

\section[tamponare]{tamponare — to rear-end, to bump into (the car in front)}

\label{lesson:IT-C413-tamponare}



\begin{quote}
Before the new one: say the Italian for to run over, then the Italian for the fine.
\end{quote}

\subsection*{You'll want to know: tamponare}
\textbf{tamponare} — \textquotedblleft{}to rear-end, to bump into (the car in front)\textquotedblright{}.

\begin{quote}
Un furgone ci ha tamponato al semaforo.
\end{quote}

\begin{quote}
\emph{A van ran into the back of us at the lights.}
\end{quote}

Also to stop a leak or a wound: \emph{tamponare una ferita}.

\begin{cousinweb}
\emph{tamponare} is built on \emph{tampone} 'plug, pad', which was borrowed from French \emph{tampon}.
\end{cousinweb}

\begin{grammarlens}[title={What you've built}]
One more word for telling, explaining and describing at B1.
\end{grammarlens}

\subsection*{Your turn}
\begin{itemize}
\raggedright
  \item \emph{Say it:} \emph{tamponare}
  \item \emph{Say it:} \emph{tamponare}, once more
  \item \emph{Say it:} a sentence of your own with \emph{tamponare}
  \item \emph{Recall:} say the Italian for to run over, then the Italian for the fine, then say \emph{tamponare} again
\end{itemize}

\begin{tcolorbox}[breakable,colback=teal!4,colframe=teal!35!black,title={Before you move on}]
What does tamponare mean? (\textquotedblleft{}To rear-end, to bump into (the car in front)\textquotedblright{}.) Say it once more.
\end{tcolorbox}
\section[il tamponamento]{il tamponamento — the rear-end collision, the pile-up}

\label{lesson:IT-C413-tamponamento}



\begin{quote}
Before the new one: say the Italian for the fine, then the Italian for to rear-end.
\end{quote}

\subsection*{You'll want to know: il tamponamento}
\textbf{il tamponamento} — \textquotedblleft{}the rear-end collision, the pile-up\textquotedblright{}.

\begin{quote}
C'è stato un tamponamento in tangenziale.
\end{quote}

\begin{quote}
\emph{There was a pile-up on the ring road.}
\end{quote}

\emph{Un tamponamento a catena} — a multi-car pile-up.

\begin{cousinweb}
\emph{tamponamento} is built on \emph{tamponare} 'to bump into', from \emph{tampone} 'plug', which was borrowed from French \emph{tampon}.
\end{cousinweb}

\begin{grammarlens}[title={What you've built}]
One more word for telling, explaining and describing at B1.
\end{grammarlens}

\subsection*{Your turn}
\begin{itemize}
\raggedright
  \item \emph{Say it:} \emph{il tamponamento}
  \item \emph{Say it:} \emph{il tamponamento}, once more
  \item \emph{Say it:} a sentence of your own with \emph{il tamponamento}
  \item \emph{Recall:} say the Italian for the fine, then the Italian for to rear-end, then say \emph{il tamponamento} again
\end{itemize}

\begin{tcolorbox}[breakable,colback=teal!4,colframe=teal!35!black,title={Before you move on}]
What does il tamponamento mean? (\textquotedblleft{}The rear-end collision, the pile-up\textquotedblright{}.) Say it once more.
\end{tcolorbox}
\section[l'ingorgo]{l'ingorgo — the traffic jam}

\label{lesson:IT-C413-ingorgo}



\begin{quote}
Before the new one: say the Italian for to rear-end, then the Italian for the rear-end collision.
\end{quote}

\subsection*{You'll want to know: l'ingorgo}
\textbf{l'ingorgo} — \textquotedblleft{}the traffic jam\textquotedblright{}.

\begin{quote}
Siamo rimasti in un ingorgo per un'ora.
\end{quote}

\begin{quote}
\emph{We were stuck in a traffic jam for an hour.}
\end{quote}

Worse than \emph{la coda}: nothing moves.

\begin{cousinweb}
\emph{ingorgo} is built from \emph{in-} and \emph{gorgo} 'whirlpool', which is inherited from Latin \emph{gurges} 'whirlpool' (EtymDB): traffic swirling and stuck.
\end{cousinweb}

\begin{grammarlens}[title={What you've built}]
One more word for telling, explaining and describing at B1.
\end{grammarlens}

\subsection*{Your turn}
\begin{itemize}
\raggedright
  \item \emph{Say it:} \emph{l'ingorgo}
  \item \emph{Say it:} \emph{l'ingorgo}, once more
  \item \emph{Say it:} a sentence of your own with \emph{l'ingorgo}
  \item \emph{Recall:} say the Italian for to rear-end, then the Italian for the rear-end collision, then say \emph{l'ingorgo} again
\end{itemize}

\begin{tcolorbox}[breakable,colback=teal!4,colframe=teal!35!black,title={Before you move on}]
What does l'ingorgo mean? (\textquotedblleft{}The traffic jam\textquotedblright{}.) Say it once more.
\end{tcolorbox}
\section[la deviazione]{la deviazione — the diversion, the detour}

\label{lesson:IT-C413-deviazione}



\begin{quote}
Before the new one: say the Italian for the rear-end collision, then the Italian for the traffic jam.
\end{quote}

\subsection*{You'll want to know: la deviazione}
\textbf{la deviazione} — \textquotedblleft{}the diversion, the detour\textquotedblright{}.

\begin{quote}
Per i lavori c'è una deviazione.
\end{quote}

\begin{quote}
\emph{Because of the roadworks there is a diversion.}
\end{quote}

\emph{Seguire la deviazione} — to follow the diversion.

\begin{cousinweb}
\emph{deviazione} is a learned borrowing from Late Latin \emph{dēviātiō}, from \emph{dēviāre} 'to turn aside', built on \emph{via} 'way' (Lewis \& Short).
\end{cousinweb}

\begin{grammarlens}[title={What you've built}]
One more word for telling, explaining and describing at B1.
\end{grammarlens}

\subsection*{Your turn}
\begin{itemize}
\raggedright
  \item \emph{Say it:} \emph{la deviazione}
  \item \emph{Say it:} \emph{la deviazione}, once more
  \item \emph{Say it:} a sentence of your own with \emph{la deviazione}
  \item \emph{Recall:} say the Italian for the rear-end collision, then the Italian for the traffic jam, then say \emph{la deviazione} again
\end{itemize}

\begin{tcolorbox}[breakable,colback=teal!4,colframe=teal!35!black,title={Before you move on}]
What does la deviazione mean? (\textquotedblleft{}The diversion, the detour\textquotedblright{}.) Say it once more.
\end{tcolorbox}
\section[deviare]{deviare — to divert, to turn off}

\label{lesson:IT-C413-deviare}



\begin{quote}
Before the new one: say the Italian for the traffic jam, then the Italian for the diversion.
\end{quote}

\subsection*{You'll want to know: deviare}
\textbf{deviare} — \textquotedblleft{}to divert, to turn off\textquotedblright{}.

\begin{quote}
Hanno deviato il traffico verso il centro.
\end{quote}

\begin{quote}
\emph{They diverted the traffic towards the centre.}
\end{quote}

Also without an object: \emph{la strada devia a destra}.

\begin{cousinweb}
\emph{deviare} is a learned borrowing from Late Latin \emph{dēviāre} 'to turn aside', from \emph{dē} 'off' and \emph{via} 'way' (Lewis \& Short).
\end{cousinweb}

\begin{grammarlens}[title={What you've built}]
One more word for telling, explaining and describing at B1.
\end{grammarlens}

\subsection*{Your turn}
\begin{itemize}
\raggedright
  \item \emph{Say it:} \emph{deviare}
  \item \emph{Say it:} \emph{deviare}, once more
  \item \emph{Say it:} a sentence of your own with \emph{deviare}
  \item \emph{Recall:} say the Italian for the traffic jam, then the Italian for the diversion, then say \emph{deviare} again
\end{itemize}

\begin{tcolorbox}[breakable,colback=teal!4,colframe=teal!35!black,title={Before you move on}]
What does deviare mean? (\textquotedblleft{}To divert, to turn off\textquotedblright{}.) Say it once more.
\end{tcolorbox}
